\pagecolor{PageColor}
\begin{center}
  {\fontsize{30pt}{26pt}\selectfont 各地区数学大学入学考试难度对比
  }
\end{center}
\separator
\begin{enumerate}
\item 印度 JEE Advanced \worldflag{IN} %add jee main
\begin{itemize}
\item 分两场（Paper 1 \& 2）,每场 180 分钟。全客观题,但含单选、多选、数值填空及矩阵匹配。
\item 特点:不仅是难,更是“阴险”。多选题选错会倒扣分（Negative Marking）,逼着你在考场上做博弈。题目不按常理出牌,常在微积分和复数里埋伏现代数学的影子,是纯粹的智力与心脏抗压测试。
\end{itemize}
\item 中国高考 (新课标 I 卷) \worldflag{CN}
\begin{itemize}
    \item 120 分钟。8 单选 + 3 多选 + 3 填空 + 5 大题。
    \item 在2021年高考上总共有8套试卷，分别为全国甲卷、全国乙卷、新高考Ⅰ卷、新高考Ⅱ卷，北京、天津、上海、浙江自主命制4套。

全国甲卷：云南、广西、贵州、四川、西藏 

全国乙卷：河南、山西、江西、安徽、甘肃、青海、内蒙古、黑龙江、吉林、宁夏、新疆、陕西 ；

新高考Ⅰ卷：山东、福建、湖北、江苏、广东、湖南、河北

新高考Ⅱ卷：海南、辽宁、重庆
    \item 特点:现在的趋势是“反套路”。虽然大纲没变,但解析几何和导数的计算量大得惊人,压轴题越来越像数学竞赛初赛。这不仅考脑力,更是在压榨考生的极限算力和在繁琐步骤中不出错的稳定性。
\end{itemize}
\item 越南高考 (THPT) \worldflag{VN}
\begin{itemize}
    \item 90 分钟。50 道单项选择题。
    \item 特点:暴力计算的巅峰。平均 1.8 分钟处理一题,基本没有喘息机会。题目充斥着大量复杂的函数套函数、立体几何最值,考生必须像人肉计算机一样扫视题目,对套路极其熟练才能做完。
\end{itemize}
\item 马来西亚统考 (UEC) \worldflag{MY}
\begin{itemize}
    \item 高级数学(I) 120 分钟；高级数学(II) 180 分钟。选择题结合大题。
    \item 特点:风格非常“老派”且扎实。不玩花哨的,就是硬碰硬的代数和证明。范围极广,连一阶微分方程和棣莫弗定理都敢考,非常看重考生的数学基本功和写完整解题步骤的习惯。
\end{itemize}
\item 韩国 CSAT (修能) \worldflag{KR}
\begin{itemize}
    \item 100 分钟。30 道填空与选择题。
    \item 特点:所谓的“杀手题”策略。前面的题都是送分,但最后两三道题（如第 22 题、30 题）的难度陡增,逻辑弯路极多。这种考试容错率低得吓人,名校生基本默认前 27 题全对,生死全看最后那几道逻辑断层题。
\end{itemize}
\item 日本 东大/京大二次试验 \worldflag{JP}
\begin{itemize}
    \item 150 分钟。6 道长篇证明/叙述题。
    \item 特点:数学界的“申论”。没有选择题,给你白纸一张让你从头推到尾。它不考反应速度,考的是你能不能把一个深刻的数论或组合问题讲清楚。题目非常优雅,是亚洲最接近“数学研究”状态的考试。
\end{itemize}
\item 台湾 分科测验 (数学甲) \worldflag{TW}
\begin{itemize}
    \item 80 分钟。单选、多选、选填及混合题（含非选择题）。
    \item 特点:短小精悍,侧重定义。由于时间只有 80 分钟,它更看重你对微积分、空间向量这些概念的直观理解。题目灵活,喜欢跨章节组合,如果概念理解不透,光靠刷题很难拿高分。
\end{itemize}
\item 香港 DSE (M1/M2) \worldflag{HK}
\begin{itemize}
    \item 150 分钟。结构化题目（Section A \& B）。
    \item 特点:职场式的模块化考核。M1 侧重应用,M2 侧重理论,题型极其固定且结构清晰。只要把往年真题（Past Paper）练透,基本能精准预测考点。它考查的是你的熟练度以及对标准化流程的执行力。
\end{itemize}
\end{enumerate}
\begin{table}[htbp]
    \centering
    \small
    \caption{全球主要数学大学入学考试：基于选拔压力的多维对比}
    \vspace{2mm}
    \begin{tabularx}{\textwidth}{lXXXXXXXX}
        \toprule
        考试项目 & 广度 & 深度 & 巧度 & 创新 & 节奏 & 负荷 & 交融 & 容错 \\ 
        \midrule
        印度 JEE Advanced & \rating{5} & \rating{5} & \rating{5} & \rating{5} & \rating{4} & \rating{4} & \rating{5} & \rating{1} \\ 
        中国高考 (新课标I) & \rating{3} & \rating{4} & \rating{4} & \rating{4} & \rating{4} & \rating{5} & \rating{5} & \rating{1} \\ 
        日本 东大/京大二次 & \rating{4} & \rating{5} & \rating{4} & \rating{5} & \rating{2} & \rating{5} & \rating{4} & \rating{2} \\ 
        越南 HGE (高中毕业) & \rating{3} & \rating{3} & \rating{3} & \rating{2} & \rating{5} & \rating{5} & \rating{2} & \rating{1} \\ 
        韩国 CSAT (修能) & \rating{3} & \rating{3} & \rating{5} & \rating{4} & \rating{5} & \rating{4} & \rating{4} & \rating{1} \\ 
        台湾 分科测验 & \rating{4} & \rating{4} & \rating{3} & \rating{3} & \rating{3} & \rating{3} & \rating{3} & \rating{3} \\ 
        香港 DSE (M1/M2) & \rating{4} & \rating{4} & \rating{2} & \rating{2} & \rating{3} & \rating{3} & \rating{3} & \rating{4} \\ 
        马来西亚统考 (UEC) & \rating{3} & \rating{3} & \rating{2} & \rating{2} & \rating{2} & \rating{4} & \rating{2} & \rating{4} \\ 
        新加坡 A-Level (H2/H3) & \rating{5} & \rating{4} & \rating{3} & \rating{3} & \rating{3} & \rating{3} & \rating{4} & \rating{2} \\ 
        马来西亚 STPM & \rating{4} & \rating{4} & \rating{3} & \rating{2} & \rating{3} & \rating{4} & \rating{3} & \rating{2} \\ 
        \bottomrule
    \end{tabularx}
\end{table}

\begin{itemize}
    \item \textbf{广度}：考察知识点的覆盖面及与大学数学（如线性代数、复分析）的衔接程度。
    \item \textbf{深度}：逻辑推理链条的长度、抽象思维能力以及对数学本质的理解。
    \item \textbf{巧度}：思维陷阱、非标准设问路径的隐蔽性及考场心理干扰。
    \item \textbf{创新}：原创题型出现频率、新定义构造及对套路化解法的破除力度。
    \item \textbf{节奏}：单题分配时间极短带来的限时压力与熟练度要求。
    \item \textbf{负荷}：突破解题切入点后，后续代数变形与纯计算量的消耗程度。
    \item \textbf{交融}：跨模块知识点的综合融合能力，要求在单题中贯通几何、代数、概率等不同领域。
    \item \textbf{容错}：得分机制的容错空间；评分越低表示顶级名校选拔时对失误近乎“零容忍”。
\end{itemize}

\begin{itemize}
    \item \textbf{Scope}: Syllabus breadth and exposure to university-level bridging topics.
    \item \textbf{Depth}: Length of logical reasoning chains and structural abstraction.
    \item \textbf{Trickiness}: Cognitive traps, non-standard setups, and psychological pressure.
    \item \textbf{Innovation}: Unseen problem mechanics, non-routine questions, and novel mathematical definitions.
    \item \textbf{Pacing}: Time pressure per question and required solution fluency.
    \item \textbf{Load}: Raw algebraic and computational intensity after finding the core approach.
    \item \textbf{Intertwinability}: Cross-domain integration combining multiple distinct math branches within a single problem.
    \item \textbf{Error Margin}: Forgiveness in scoring; a lower rating indicates near-zero tolerance for mistakes for elite admissions.
\end{itemize}

\centering
% Usage: \radarChart{Exam Title}{Color}{8 Scores in Order}
\radarChart{India JEE Advanced}{jeeRed}{5, 5, 5, 5, 4, 4, 5, 1}

\vspace{1cm}

\radarChart{China Gaokao (New I)}{gaokaoBlue}{3, 4, 4, 4, 4, 5, 5, 1}